% Generated from locale/tr/content/set-theory/choice/vitali.tex; do not hand-edit.


\olfileid[tr]{sth}{choice}{vitali}
\olsection{Ek: Vitali Paradoksu}

Banach--Tarski kuruluşunun kabul edilebilir olup olmadığı konusunda gerçek
bir fikir edinmek için \emph{ispatını} incelemeliyiz. Ne yazık ki bu,
burada sunabileceğimizden çok daha fazla cebir gerektirir. Bununla birlikte
şu sonuca odaklanarak Banach--Tarski ispatına biraz ışık tutabilecek kısa
açıklamalar sunabiliriz:\footnote{Çok daha kapsamlı bir inceleme için
\cite{Weston2003} ya da \cite{Wagon2016} kaynaklarına bakın.}

\begin{thm}[Vitali Paradoksu ($\ZFC$ içinde)]\ollabel{vitaliparadox}
Her çember sayılabilir çoklukta parçaya ayrılabilir ve bu parçalar
(döndürme ve öteleme yoluyla) yeniden bir araya getirilerek o çemberin iki
kopyası oluşturulabilir.
\end{thm}

Vitali Paradoksu'nu ispatlamak Banach--Tarski Paradoksu'nu ispatlamaktan
çok daha kolaydır. Ona ``Vitali Paradoksu'' dedik; çünkü Vitali'nin
ölçülemeyen bir küme veren \citeyear{Vitali1905} tarihli kuruluşundan
çıkar. Fakat Vitali Paradoksu ile Banach--Tarski Paradoksu ispatlarının
küme teorisiyle ilgili yönleri birbirine çok benzer. Sonuçlar arasındaki
temel fark yalnızca şudur: Banach--Tarski \emph{sonlu} bir ayrıştırmayı,
Vitali Paradoksu ise \emph{sayılabilir sonsuz} bir ayrıştırmayı ele alır.
\citet{Weston2003} şöyle der: Vitali Paradoksu ``kesinlikle Banach--Tarski
paradoksu kadar çarpıcı değildir, fakat geometrik paradoksların `basit'
durumlarda bile ortaya çıkabileceğini gösterir.''

Vitali Paradoksu iki boyutlu bir şekil olan çemberle ilgilidir. Dolayısıyla
$\Real^2$ düzleminde çalışacağız. $\rotationsgroup$, noktaları başlangıç
noktası çevresinde $[0,2\pi)$ arasındaki \emph{rasyonel} radyan değerleri
kadar (saat yönünde) döndüren dönüşümlerin kümesi olsun. İşte
$\rotationsgroup$ hakkında bazı cebirsel olgular (sonucun ifadesini
anlamazsanız ispat anlamını açıklayacaktır):

\begin{lem}\ollabel{rotationsgroupabelian}
$\rotationsgroup$, fonksiyon bileşkesi altında değişmeli bir {grup}
oluşturur.
\end{lem}

\begin{proof}
$0$ radyanlık döndürme için $0_{\rotationsgroup}$ yazalım. Bu,
$\rotationsgroup$ için birim elemandır; çünkü her
$\rho\in\rotationsgroup$ için
$\comp{0_{\rotationsgroup}}{\rho}=\comp{\rho}{0_{\rotationsgroup}}=
\rho$ olur.

Her elemanın bir tersi vardır. $\rho\in\rotationsgroup$, $r$ radyan
döndürüyorsa $r=0$ olduğunda $\rho^{-1}$ de $0$ radyan, aksi hâlde
$2\pi-r$ radyan döndürür; böylece $\rho^{-1}\in\rotationsgroup$ ve
$\rho\circ\rho^{-1}=0_\rotationsgroup$ olur.

Bileşke birleşmelidir: her
$\rho,\sigma,\tau\in\rotationsgroup$ için
$\comp{\rho}{(\comp{\sigma}{\tau})}=
\comp{(\comp{\rho}{\sigma})}{\tau}$.

Bileşke değişmelidir: her $\rho,\sigma\in\rotationsgroup$ için
$\comp{\rho}{\sigma}=\comp{\sigma}{\rho}$.
\end{proof}

Gerçekten $\rotationsgroup$ grubumuzu iki yarıya ayırıp bütün grubu geri
elde etmek için bu yarıların herhangi birini kullanabiliriz:

\begin{lem}\ollabel{disjointgroup}
$\rotationsgroup$ için her ikisi de birer taban olan iki ayrık
$\rotationsgroup_{1}$ ve $\rotationsgroup_{2}$ kümesine ayrılmış bir
parçalanış vardır.
\end{lem}

\begin{proof}
$\rotationsgroup_{1}$, $[0,\pi)$ aralığındaki rasyonel radyan değerleri
kadar döndürmelerden oluşsun;
$\rotationsgroup_2=\rotationsgroup\setminus\rotationsgroup_1$ olsun.
Temel cebir gereği
$\Setabs{\comp{\rho}{\rho}}{\rho\in\rotationsgroup_1}=
\rotationsgroup$. $\rotationsgroup_2$ için de benzer bir sonuç elde edilir.
\end{proof}

Gruplar hakkındaki bu olguyu \olref{vitaliparadox} sonucunu kurmak için
kullanacağız. $\onesphere$ birim çember, yani düzlemin başlangıç noktasına
tam olarak $1$ birim uzaklıktaki noktalar kümesi, başka bir deyişle
$\Setabs{\tuple{r,s}\in\Real^2}{\sqrt{r^2+s^2}=1}$ olsun. $\onesphere$
üzerinde şu bağıntıyı ele alarak $\onesphere$'ı parçalara ayıracağız:
\[
	r \sim s \emph{ ancak ve ancak }(\exists \rho \in \rotationsgroup)\rho(r) = s.
\]
Yani $\onesphere$'ın iki noktası, ancak ve ancak başlangıç noktası
çevresinde rasyonel değerli bir döndürmeyle birinden ötekine
ulaşabiliyorsak bu bağıntıyla bağlantılıdır. Şaşırtıcı olmayan biçimde:

\begin{lem}
$\sim$ bir denklik bağıntısıdır.
\end{lem}

\begin{proof}
\olref{rotationsgroupabelian} kullanılarak sıradandır.
\end{proof}

Şimdi Seçim'i kullanarak $\onesphere$'ın $\sim$ altındaki her denklik
sınıfından tam bir eleman içeren bir $C$ kümesi elde ederiz. Yani denklik
sınıfları kümesi üzerinde bir $f$ seçim fonksiyonu ele alırız,\footnote{
$\rotationsgroup$ !!{enumerable} olduğundan $E$'nin her !!{element}
!!{enumerable}. $\onesphere$ !!{nonenumerable} olduğundan
\olref{vitalicover} ve
\olref[card-arithmetic][simp]{kappaunionkappasize} sonuçları, $E$'nin
!!{nonenumerable} olduğunu gerektirir. Dolayısıyla bu, \emph{sayılamaz}
Seçim'in bir kullanımıdır.}
\[
	E = \Setabs{\equivrep{r}{\sim}}{r \in \onesphere},
\]
ve $C=\ran{f}$ olsun. Her $\rho\in\rotationsgroup$ döndürmesi için
$\funimage{\rho}{C}$ kümesi, $\rho$ döndürmesinin $C$'deki her noktaya
uygulanmasıyla elde edilen noktalardan oluşur. Sıradaki iki sonuç bu
kümelerin çemberi bütünüyle ve örtüşmeden kapladığını gösterir:

\begin{lem}\ollabel{vitalicover}
$\onesphere = \bigcup_{\rho \in \rotationsgroup} \funimage{\rho}{C}$.
\end{lem}

\begin{proof}
$s\in\onesphere$ sabit olsun. $r\in\equivrep{s}{\sim}$, yani
$r\sim s$, yani bir $\rho\in\rotationsgroup$ için $\rho(r)=s$ olan bir
$r\in C$ vardır.
\end{proof}

\begin{lem}\ollabel{vitalinooverlap}
$\rho_1\neq\rho_2$ ise
$\funimage{\rho_1}{C}\cap\funimage{\rho_2}{C}=\emptyset$.
\end{lem}

\begin{proof}
$s\in\funimage{\rho_1}{C}\cap\funimage{\rho_2}{C}$ varsayın. Öyleyse
bir $r_1,r_2\in C$ için
$s=\rho_1(r_1)=\rho_2(r_2)$ olur. Dolayısıyla
$\rho^{-1}_2(\rho_1(r_1))=r_2$ ve
$\comp{\rho_1}{\rho^{-1}_2}\in\rotationsgroup$; böylece
$r_1\sim r_2$ olur. $C$, $\sim$ altındaki her denklik sınıfından tam bir
eleman seçtiğinden $r_1=r_2$. Dolayısıyla
$s=\rho_1(r_1)=\rho_2(r_1)$ ve böylece $\rho_1=\rho_2$.
\end{proof}

Şimdi önceki cebirsel olgularımızı çemberimize uyguluyoruz:

\begin{lem}\ollabel{pseudobanachtarski}
$\onesphere$'ın iki ayrık $D_1$ ve $D_2$ kümesine ayrılmış öyle bir
parçalanışı vardır ki $D_1$, döndürülerek $\onesphere$'ın bir kopyasını
oluşturabilecek sayılabilir çoklukta kümeye parçalanabilir (ve aynı şey
$D_2$ için de geçerlidir).
\end{lem}

\begin{proof}
\olref{disjointgroup} içindeki $\rotationsgroup_1$ ve
$\rotationsgroup_2$ kümelerini kullanarak şunları tanımlayın:
\begin{align*}
	D_{1} &= \bigcup_{\rho \in \rotationsgroup_1} \funimage{\rho}{C} &
	D_{2} &= \bigcup_{\rho \in \rotationsgroup_2} \funimage{\rho}{C}
\end{align*}
Bu, \olref{vitalicover} gereği $\onesphere$'ın bir parçalanışıdır ve
\olref{vitalinooverlap} gereği $D_1$ ile $D_2$ ayrıktır. Kuruluş gereği
$D_1$, her $\rho\in\rotationsgroup_1$ için $\funimage{\rho}{C}$ olmak
üzere sayılabilir çoklukta kümeye parçalanabilir. Bunlar döndürülerek
$\onesphere$'ın bir kopyasını oluşturabilir; çünkü
\olref{disjointgroup} ve \olref{vitalicover} gereği
$\onesphere=\bigcup_{\rho\in\rotationsgroup}\funimage{\rho}{C}=
\bigcup_{\rho\in\rotationsgroup_1}
\funimage{(\comp{\rho}{\rho})}{C}$ olur. Aynı uslamlama $D_2$'ye de
uygulanır.
\end{proof}\noindent
Bu, Vitali Paradoksu'nu doğrudan gerektirir. Çünkü $\onesphere$'ı
sayılabilir çoklukta parçaya (çeşitli $\funimage{\rho}{C}$ kümelerine)
ayırıp bunları rijit biçimde hareket ettirerek ($D_1$'in her parçasını
yalnızca döndürerek, $D_2$'nin her parçasını ise önce öteleyip sonra
döndürerek) $\onesphere$'dan \emph{iki} $\onesphere$ kopyası üretebiliriz.

İspat stratejisini özetleyelim. Düzlemdeki döndürmeler grubu hakkında bazı
cebirsel olgularla başladık. Bu grubu $\onesphere$'ı denklik sınıflarına
ayırmak için kullandık. Sonra her sınıftan elemanlar seçmek için Seçim'i
kullanarak bir ``paradoksa'' ulaştık.

Banach--Tarski'yi ispatlamak için tam olarak aynı stratejiyi kullanırız.
Temel fark, Banach--Tarski'yi ispatlamak için kullanılan cebirsel olguların
Vitali Paradoksu'nu ispatlamak için kullanılanlardan önemli ölçüde daha
karmaşık olmasıdır. Fakat bu cebirsel olguların Seçim'le hiçbir ilgisi
yoktur. Bunları kısaca özetleyeceğiz.

Banach--Tarski'yi ispatlamak için \olref{disjointgroup} sonucunun bir
benzerini kurarak başlarız: her \emph{serbest grup}, sezgisel olarak
``hareket ettirerek'' bütün grubun iki kopyasını geri elde edebileceğimiz
dört parçaya ayrılabilir.\footnote{\emph{Dört} parça kullanabilmemiz
\cite{Robinson1947} sonucudur. Güncel bir ispat için
\citet[Theorem 5.2]{Wagon2016} eserine bakın. Grubun parçalarını
``hareket ettirme'' betimlemesinde \citet[p.~3]{Weston2003} eserini
izliyoruz.} Sonra $\Real^3$'ün başlangıç noktası çevresindeki iki belirli
döndürmeyi kullanarak serbest bir döndürmeler grubu $F$ üretebileceğimizi
gösteririz.\footnote{Bkz. \citet[Theorem 2.1]{Wagon2016}.} (Henüz Seçim
yok.) Şimdi kürenin yüzeyindeki noktaları, ancak ve ancak biri ötekinden
$F$ içindeki bir döndürmeyle elde edilebiliyorsa ``benzer'' sayarız. Sonra
``benzer'' noktaların her denklik sınıfından tam bir nokta seçmek için
\emph{Seçim'i kullanırız}. $F$'yi bölüşümüzü kürenin yüzeyine
\olref{pseudobanachtarski} içindeki gibi uygulayıp yüzeyi dört parçaya
ayırır ve bunları ``hareket ettirerek'' küre yüzeyinin iki kopyasını elde
ederiz. Böylece \citep{Hausdorff1914} kurulur:

\begin{thm}[Hausdorff Paradoksu ($\ZFC$ içinde)]
Her kürenin yüzeyi sonlu sayıda parçaya ayrılabilir ve bu parçalar
(döndürme ve öteleme yoluyla) yeniden bir araya getirilerek o kürenin iki
ayrık kopyası oluşturulabilir.
\end{thm}

Tam Banach--Tarski Teoremi'ni (yalnızca kürenin yüzeyiyle değil, içiyle de
ilgilenen teoremi) elde etmek için birkaç ek cebirsel hile gerekir. Fakat
açıkçası bu yalnızca cebir pastasının süsüdür. Weston bu nedenle şöyle
yazar:
\begin{quote}
  [\ldots] serbest gruplar hakkındaki sonuç, Banach--Tarski paradoksunun
  ispatındaki \emph{temel adımdır}. Bu bakış açısından Banach--Tarski
  paradoksu, $\Real^3$ hakkındaki bir ifade olmaktan çok [$\Real^3$'teki
  öteleme ve döndürmeler] grubunun karmaşıklığı hakkındaki bir ifadedir.
  \cite[p.~16]{Weston2003}
\end{quote}
Yani \emph{sonlu} bir ayrıştırma (Banach--Tarski'deki gibi) mı yoksa
\emph{sayılabilir sonsuz} bir ayrıştırma (Vitali Paradoksu'ndaki gibi) mı
sunabileceğimiz, iki boyutta ya da üç boyutta çalışmak hakkındaki belirli
grup kuramsal olgulara dayanır.

Kabul etmek gerekir ki bu son gözlem \olref[banach]{sec} sonundaki şakayı
biraz bozar. ``Banach--Tarski'' iki boyutlu olduğundan parçalarını yeniden
düzenleyip ``Banach--Tarski Banach--Tarski'' oluşturmak istiyorsak
sayılabilir \emph{sonsuzlukta} parçaya ayrılmalıdır. Şakayı düzeltmek için
üç boyutta yazmak gerekir. Bunu okura alıştırma olarak bırakıyoruz.

Son bir yorum. \olref[banach]{sec} içinde elde edilen küre ``parçalarının''
\emph{ölçülebilir} olamayacağını, gözde canlandırılamayan ``sonsuz
saçılımlar'' olmaları gerektiğini belirttik. Seçim'i
\olref{pseudobanachtarski} sonucunu elde etmek için kullanışımızda da aynı
şey geçerlidir. Bütün bunlar açıklanmaya değer.

Yine biraz arka planın taslağını çizmeliyiz (fakat bu \emph{yalnızca} bir
taslaktır; \emph{ölçü} hakkındaki bir ders kitabı maddesine başvurmak
isteyebilirsiniz). Bir $X$ kümesi için ölçü tanımlamak, $X$ üzerindeki bir
``$\sigma$-cebiri'' içinde bulunan her $E$ için bir
$\mu(E)\in\Real$ değeri atamaktır. Ayrıntılar önemli değildir; şu kadar ki
$\mu$ fonksiyonu sayılabilir toplamsallık ilkesine uymalıdır: ayrık
kümelerin sayılabilir bir birleşiminin ölçüsü onların ayrı ayrı ölçülerinin
toplamıdır; yani $X_n$'ler ayrık olduğunda
$\mu(\bigcup_{n < \omega} X_n) = \sum_{n < \omega}\mu(X_n)$ olur. Bir
kümenin ``ölçülemez'' olduğunu söylemek, ona uygun biçimde hiçbir ölçü
atanamayacağını söylemektir. Şimdi önceki $\rotationsgroup$ grubumuzu
kullanarak:

\begin{cor}[Vitali]
$\mu(\onesphere)=1$ olan ve $X$ ile $Y$ eşse
$\mu(X)=\mu(Y)$ koşulunu sağlayan bir $\mu$ ölçüsü olsun. O zaman her
$\rho\in\rotationsgroup$ için $\funimage{\rho}{C}$ ölçülemezdir.
\end{cor}

\begin{proof}
Olmayana ergiyle tersini varsayın. Öyleyse bir
$\sigma\in\rotationsgroup$ ve bir $r\in\Real$ için
$\mu(\funimage{\sigma}{C})=r$ olsun. Her
$\rho\in\rotationsgroup$ için $\funimage{\rho}{C}$ ve
$\funimage{\sigma}{C}$ eştir; dolayısıyla her
$\rho\in\rotationsgroup$ için $\mu(\funimage{\rho}{C})=r$ olur.
\olref{vitalicover} ve \olref{vitalinooverlap} gereği
$\onesphere=\bigcup_{\rho\in\rotationsgroup}\funimage{\rho}{C}$,
ikişer ikişer ayrık kümelerin sayılabilir bir birleşimidir. Öyleyse
sayılabilir toplamsallık, $\mu(\onesphere)=1$ değerinin her
$\funimage{\rho}{C}$ kümesinin ölçülerinin toplamı olmasını gerektirir;
yani:
\[
	1 = \mu(\onesphere) = \sum_{\rho \in \rotationsgroup}\mu(\funimage{\rho}{C}) = \sum_{\rho \in \rotationsgroup}r
\]
Fakat $r=0$ ise $\sum_{\rho\in\rotationsgroup}r=0$ ve $r>0$ ise
$\sum_{\rho\in\rotationsgroup}r=\infty$ olur.
\end{proof}

